% page 545 du fac-simile (image p-544 du scan)
% bloc 170.2 x 237.7 mm ; marge gauche 31.8 mm
\pgc{10.160mm}{0.000mm}{
\l{\cel{57}537}
\l{}
\l{\sou{skopo}. I. To quon onu esforcas atingar. - II. To a quo tendencas}
\l{\cel{3}ago. - III.To quon intencas kauzo inteligenta. - EI.}
\l{}
\l{\sou{skopso}. (\sou{zool.}) Genero de rapt-ucelo noktala, ek la familio \textquotedbl{}bubo-}
\l{\cel{3}nidi\textquotedbl{}. L.\cel{1}\sou{scops aldrovandi}. - EFL.}
\l{}
\l{\sou{skorbuto}. (\sou{patol.}) Morbo quan karakterizas la sango-povresko,}
\l{\cel{3}ulceri, makuli hemoragiala, la altereso di la jenjivi, la kario}
\l{\cel{3}di la denti e di la osti maxilala, abateso, deperiso, produktita}
\l{\cel{3}(normale) da la manjo tro ofta ed exkluziva, di salizaji, da la}
\l{\cel{3}privaceso di legumi fresha e vitaminoza. - DEFIRS.}
\l{}
\l{\sou{skorio}. Substanco vitratra qua acensas a la surfaco di la metali}
\l{\cel{3}fuzanta, e konsistas ye materii altra-sorta, o ye partikuli de}
\l{\cel{3}metalo. - EFIS.}
\l{}
\l{\sou{skorpiono}. (\sou{zool.}) Animalo ek la klaso \textquotedbl{}araknidi pulmonoza\textquotedbl{}, di qua}
\l{\cel{3}la kaudo havas dardo qua komunikas kun glando venenifanta.-}
\l{\cel{3}L.\cel{1}\sou{euscorpius italicus}. - (\sou{astron.})\cel{2}La signo okesma di zodiako.}
\l{\cel{3}- DEFIRS.}
\l{}
\l{\sou{skorzonero}. (\sou{bot.}) Genero di planto kompozaja, cikoriacea, quan}
\l{\cel{3}karakterizas kapituli flava o purpurea, cirkondata da braktei}
\l{\cel{3}dispozita plura-range, e frukti granda qui havas penacheto (sen-}
\l{\cel{3}stipita) de pili plumatra. - L.\cel{1}\sou{scorzonera}. DEFIRS.}
\l{}
\l{\sou{skotisho}. Danso di qua la ritmo kelke similesas olta di polko, ma}
\l{\cel{3}du-tempa (vice 2-4) e plu lenta. -\cel{2}DEFIS.}
\l{}
\l{\sou{skovelo}. I. Sorto di balayilo, ek linjo-peco ligita ye un ek la}
\l{\cel{3}extremaji di longa bastono, por netigar la forno di la pano-}
\l{\cel{3}bakeyi. - II. Sorto di balayilo o de brosilo, kun longa mancho,}
\l{\cel{3}por netigar la parto interna di kanono pos pafo. -}
\l{}
\l{skrachar. (\sou{tran}s.,\cel{1}\sou{per)}.\cel{2}Lacerar ne-profunde la pelo per la ungli}
\l{\cel{3}o per ulo pintoza. - E.}
\l{}
\l{\sou{skrapar}. (\sou{trans.)}\cel{2}Frotar til deprenar parto di la surfaco. - E.}
\l{}
\l{\sou{skreno}. I. Framo metala, kun pedi, kovrita per treliso densa,}
\l{\cel{3}quan onu lokizas avan kameno-fairuyo, por protektar su kontre}
\l{\cel{3}la ardoro tro granda di la fairo. - II. Framo ligna, kovrita}
\l{\cel{3}per telo tensita, quan la piktisto, la desegnisto, lokizas avan}
\l{\cel{3}la fenestro di sua chambro laborala por diminutar la lumo-dazlo}
\l{\cel{3}tro granda. -(\sou{cinemo}) Telo blanka, tabelo blanka, sur qua la}
\l{\cel{3}lumo-radii qui departas de filmo-projektoro, su reflektas por}
\l{\cel{3}reproduktar la imajo projektita. (analoge pri televiziono).}
\l{\cel{3}- EFR.(pri cinemo e tel.:\cel{1}\sou{+skrino})}
\l{}
\l{\sou{skribar.}\cel{2}(\sou{trans.).}\cel{1}Trasar la grupi de karakteri konvencionita por}
\l{\cel{3}reprezentar vorti. - DEFIS.}
\l{}
\l{skriptar. (netrans.) Esar verkifisto di literaturaji.}
}
